\documentclass[11pt,a4paper]{article}

\usepackage{amsmath,amssymb}
\usepackage[margin=0.9in]{geometry}
\usepackage{xcolor}
\usepackage{enumitem}
\usepackage{booktabs}
\usepackage{hyperref}
\usepackage{mdframed}

\definecolor{lightblue}{rgb}{0.90,0.95,1.0}
\definecolor{lightyellow}{rgb}{1.0,0.98,0.90}
\definecolor{lightgreen}{rgb}{0.92,0.98,0.92}
\definecolor{blueframe}{rgb}{0.4,0.6,0.8}
\definecolor{orangeframe}{rgb}{0.8,0.5,0.1}
\definecolor{greenframe}{rgb}{0.4,0.7,0.4}

\newmdenv[backgroundcolor=lightyellow,linecolor=orangeframe,linewidth=1pt,innerleftmargin=10pt,innerrightmargin=10pt,innertopmargin=6pt,innerbottommargin=6pt]{onelinerbox}
\newmdenv[backgroundcolor=lightblue,linecolor=blueframe,linewidth=1pt,innerleftmargin=10pt,innerrightmargin=10pt,innertopmargin=6pt,innerbottommargin=6pt]{whytbox}
\newmdenv[backgroundcolor=lightgreen,linecolor=greenframe,linewidth=1pt,innerleftmargin=10pt,innerrightmargin=10pt,innertopmargin=6pt,innerbottommargin=6pt]{codeboxmd}
\newmdenv[backgroundcolor=white,linecolor=gray,linewidth=1pt,innerleftmargin=10pt,innerrightmargin=10pt,innertopmargin=6pt,innerbottommargin=6pt]{stepboxmd}

\newcommand{\step}[1]{\noindent\textbf{#1}\par\smallskip}
\newcommand{\oneliner}[1]{\begin{onelinerbox}\textbf{If asked about this, say:} #1\end{onelinerbox}\vspace{6pt}}
\newcommand{\whyt}[1]{\begin{whytbox}\textbf{Why do we need this?} #1\end{whytbox}\vspace{6pt}}
\newcommand{\codebox}[1]{\begin{codeboxmd}\textbf{The code:} #1\end{codeboxmd}\vspace{6pt}}
\newcommand{\stepboxbegin}{\begin{stepboxmd}\textbf{Step-by-step}\par\smallskip}
\newcommand{\stepboxend}{\end{stepboxmd}\vspace{6pt}}

\title{Team Math Guide\\
\normalsize Cheat sheet + full derivations for DCA-Trie}
\author{Bernard}
\date{\today}

\begin{document}
\maketitle

%% ============================================================
%% PART 1: CHEAT SHEET
%% ============================================================
\newpage
\section*{Cheat Sheet --- All Key Equations at a Glance}
\addcontentsline{toc}{section}{Cheat Sheet}

\subsection*{Notation}

\begin{center}
\begin{tabular}{@{}lll@{}}
\toprule
Symbol & Meaning & Plain English \\
\midrule
$\mathcal{G}$ & Knowledge graph & A network of facts (entities + relations) \\
$e \in \mathcal{E}$ & Entity node & A named thing (e.g.\ ``Paris'', ``France'') \\
$r \in \mathcal{R}$ & Relation & A typed edge (e.g.\ ``capital\_of'') \\
$\vec{c}$ & Reasoning chain & A path through the KG: $c_1 \xrightarrow{r_1} c_2 \xrightarrow{r_2} \cdots$ \\
$L$ & Chain length & Number of hops in the path \\
$\mathcal{N}(s)$ & Neighbours of $s$ & All nodes directly connected to $s$ \\
$\mathcal{N}(s,r)$ & Neighbours via $r$ & Nodes connected to $s$ by relation $r$ \\
$L_{\text{lm}}$ & LLM & The language model we use for generation \\
$\text{head}(i)$ & Head entity at step $i$ & The entity we are currently expanding \\
$\tau$ & Cosine threshold & Minimum similarity to accept an entity \\
$\mathcal{T}$ & TypeOracle & KG-based type checker \\
$\text{rk}_e$ & KG role of entity $e$ & What types $e$ plays (e.g.\ person, city) \\
$\text{tr}_r$ & Range of relation $r$ & What type the target must be \\
$T(\vec{c})$ & Correctness gate & 1 if all triples are in the KG \\
$C_{\text{max}}$ & Max valid chains & Total number of real KG paths of length $L$ \\
$B$ & Beam size & How many paths we keep in parallel \\
$d_{\text{avg}}$ & Average out-degree & Mean number of neighbours per node \\
$S$ & Score & How good a reasoning path is \\
$S^*$ & Best score & Highest scoring valid answer \\
\bottomrule
\end{tabular}
\end{center}

\subsection*{Key Equations}

\begin{enumerate}[leftmargin=*,itemsep=6pt]

\item \textbf{Path probability (chain rule):}
\[
P(\vec{c}) = \prod_{i=1}^{L} P(c_i \mid c_1, c_2, \dots, c_{i-1})
\]

\item \textbf{Graph-constrained logits:}
\[
p_i(v) = \begin{cases}
\text{softmax}(z_i)_v & \text{if } v \in \mathcal{N}(\text{head}(i)) \\
0 & \text{otherwise}
\end{cases}
\]

\item \textbf{Cosine similarity (v1 scoring):}
\[
\text{sim}(a, b) = \frac{a \cdot b}{\|a\|\,\|b\|}
\]

\item \textbf{Decomposed product (v2 scoring):}
\[
S(\vec{c}) = \prod_{i=1}^{L} P(c_i \mid \vec{c}_{<i}) \cdot T(\vec{c})
\]
where $T(\vec{c}) = \prod_{i=1}^{L} \mathbb{1}[(\text{head}(i),\, r_i,\, c_{i+1}) \in \mathcal{G}]$.

\item \textbf{TypeOracle range gate:}
\[
\text{range\_gate}(r, e) = \mathbb{1}\big[\text{tr}_r \subseteq \text{rk}_e\big]
\]

\item \textbf{TypeOracle type gate:}
\[
\text{type\_gate}(s, r) = \bigcup_{e \in \mathcal{N}(s,r)} \text{rk}_e
\]

\item \textbf{Hits@k:}
\[
\text{Hits@}k = \frac{|\{q : \text{rank}(q) \leq k\}|}{|\mathcal{Q}|}
\]

\item \textbf{V1 complexity:}
\[
O(E^L \cdot L)
\]
where $E = \max_{s \in \mathcal{G}} |\mathcal{N}(s)|$.

\item \textbf{V2 complexity:}
\[
O(L \cdot B \cdot d_{\text{avg}})
\]

\end{enumerate}

\subsection*{One-Liners for Defense}

\begin{center}
\small
\begin{tabular}{@{}p{5.5cm}p{9cm}@{}}
\toprule
If they ask\dots & You say\dots \\
\midrule
``What is graph-constrained decoding?'' &
``We mask the LLM's output at each step so it can only generate tokens that are valid edges in the knowledge graph.'' \\[4pt]
``Why cosine similarity?'' &
``It measures how close the LLM's prediction is to the real KG embedding, without caring about magnitude---just direction.'' \\[4pt]
``What does the TypeOracle do?'' &
``It's a semantic gate: before we let the LLM generate a relation, we check whether its target type is consistent with the entity's role in the KG.'' \\[4pt]
``Why is v2 better than v1?'' &
``V1 scales exponentially with chain length $O(E^L)$. V2 scales linearly $O(L \cdot B \cdot d)$ by using beam search instead of enumerating all paths.'' \\[4pt]
``What does $T(\vec{c})$ mean?'' &
``It's a correctness check: 1 if every triple in the chain actually exists in the KG, 0 otherwise. This is what makes DCA-Trie faithful.'' \\[4pt]
``Is the output faithful?'' &
``Yes. We prove that the constrained decoder only outputs chains where every triple is in the KG. The proof uses mathematical induction on the chain length.'' \\
\bottomrule
\end{tabular}
\end{center}

%% ============================================================
%% PART 2: FULL DERIVATIONS
%% ============================================================
\newpage
\section*{Full Derivations}
\addcontentsline{toc}{section}{Full Derivations}

%% ------------------------------------------------------------
\subsection{1. Path Scoring --- Why a Product?}

\whyt{
We need to assign a score to each reasoning path $\vec{c} = (c_1, r_1, c_2, r_2, \dots, c_L)$. The score tells us how likely this path is according to the LLM. We want high-probability paths (the LLM is confident) to rank higher.
}

\stepboxbegin
\step{Start with the joint probability.}
We want $P(c_1, r_1, c_2, r_2, \dots, c_L)$---the probability of the entire sequence of tokens.

\step{Apply the chain rule of probability.}
The chain rule says: the joint probability of a sequence equals the product of each element conditional on everything before it:
\[
P(x_1, x_2, \dots, x_n) = P(x_1) \cdot P(x_2 \mid x_1) \cdot P(x_3 \mid x_1, x_2) \cdots P(x_n \mid x_1, \dots, x_{n-1})
\]
This is always exact---no approximation. It's just probability theory.

\step{Apply to our tokens.}
Each step we generate either an entity name or a relation type. So:
\[
P(\vec{c}) = \prod_{i=1}^{L} P(\text{entity}_i \mid \text{all previous tokens}) \cdot P(\text{relation}_i \mid \text{all previous tokens})
\]

\step{That's it.}
The path score is just the LLM's probability of generating that exact sequence of tokens, step by step.
\stepboxend

\codebox{
The path score is computed in \texttt{dca\_v2\_generate()} at \texttt{experiments/type\_oracle\_full/decoding.py}. It multiplies the step-by-step LLM probabilities for each candidate path.
}

\oneliner{
``The path score is just the LLM's own probability of generating that sequence---we multiply the step-by-step probabilities using the chain rule.''
}

%% ------------------------------------------------------------
\subsection{2. Graph-Constrained Logits --- How We Enforce the KG}

\whyt{
Without constraints, the LLM can generate any token in its vocabulary (50,000+ tokens). We want to restrict it to only tokens that correspond to valid KG edges from the current entity.
}

\stepboxbegin
\step{The LLM outputs logits.}
At each step, the LLM produces a vector $z \in \mathbb{R}^{|V|}$ where $|V|$ is the vocabulary size. Each entry $z_v$ is a raw score for token $v$.

\step{We look up valid tokens.}
Given the current head entity $\text{head}(i)$, we find all neighbours in the KG:
\[
\mathcal{N}(\text{head}(i)) = \{v : (\text{head}(i), r, v) \in \mathcal{G} \text{ for some } r\}
\]
These are the only entities the LLM is allowed to mention next.

\step{We mask invalid tokens.}
We set the logits of invalid tokens to $-\infty$. After softmax, these become probability 0:
\[
p_i(v) = \begin{cases}
\text{softmax}(z_i)_v & \text{if } v \in \mathcal{N}(\text{head}(i)) \\
0 & \text{otherwise}
\end{cases}
\]

\step{Result: the LLM can only generate valid KG entities.}
It's like a dropdown menu instead of a text box---you can only pick from what's actually in the graph.
\stepboxend

\codebox{
The mask is built in \texttt{GraphConstrainedDecoding} at \texttt{src/graph\_constrained\_decoding.py}. It calls \texttt{allowed\_tokens\_fn(head\_entity, allowed\_tokens\_set)} which returns the IDs of valid next tokens.
}

\oneliner{
``We set the logits of all invalid tokens to negative infinity so softmax makes their probability zero---the LLM physically cannot generate them.''
}

%% ------------------------------------------------------------
\subsection{3. Cosine Similarity --- The v1 Gate}

\whyt{
The graph constraint ensures structural validity (the triple exists). But we also want semantic validity: does the LLM \emph{mean} the same entity as the KG? Cosine similarity measures this.
}

\stepboxbegin
\step{What is cosine similarity?}
Given two vectors $a$ and $b$:
\[
\text{sim}(a, b) = \frac{a \cdot b}{\|a\|\,\|b\|} = \frac{\sum_i a_i b_i}{\sqrt{\sum_i a_i^2} \cdot \sqrt{\sum_i b_i^2}}
\]
This measures the angle between the vectors, not their length. Range: $[-1, 1]$.

\step{Why cosine, not Euclidean distance?}
Because we care about direction, not magnitude. The LLM's internal representation of ``Paris'' might have different magnitude than the KG embedding, but if they point in the same direction, they represent the same concept.

\step{How do we use it?}
At each step, we:
\begin{enumerate}
  \item Get the LLM's internal state for the entity it just generated.
  \item Get the KG embedding of the candidate entity.
  \item Compute cosine similarity.
  \item If $\text{sim} > \tau$, we accept it; otherwise we reject it.
\end{enumerate}

\step{What is $\tau$?}
The threshold. We set it using the validation set. Typical range: 0.4--0.7. Too low = accept wrong entities. Too high = reject correct ones.
\stepboxend

\oneliner{
``Cosine similarity measures the angle between the LLM's prediction and the KG embedding---same direction means same meaning, regardless of magnitude.''
}

%% ------------------------------------------------------------
\subsection{4. Decomposed Product --- The v2 Scoring Formula}

\whyt{
In v1, we scored entire paths at once. In v2, we score each step independently and multiply. This makes beam search possible and scales linearly instead of exponentially.
}

\stepboxbegin
\step{Start with the path probability.}
From Section 1:
\[
P(\vec{c}) = \prod_{i=1}^{L} P(c_i \mid c_{<i})
\]

\step{Add the correctness gate.}
We multiply by $T(\vec{c})$ which is 1 if every triple is in the KG, 0 otherwise:
\[
S(\vec{c}) = P(\vec{c}) \cdot T(\vec{c}) = \left(\prod_{i=1}^{L} P(c_i \mid c_{<i})\right) \cdot \left(\prod_{i=1}^{L} \mathbb{1}[(\text{head}(i), r_i, c_{i+1}) \in \mathcal{G}]\right)
\]

\step{Why the product of indicators?}
Each indicator checks one triple. If any triple is missing from the KG, the whole product is 0. This is how we guarantee faithfulness---a single invalid triple kills the score.

\step{Why decompose instead of scoring the whole path?}
Because we can now do beam search: at each step, keep only the top-$B$ partial paths. We don't need to enumerate all $E^L$ paths.

\step{The conditional independence assumption.}
We assume each step depends only on the entities so far, not on the full token sequence. This is approximate but practical---the LLM's hidden state captures the history.
\stepboxend

\oneliner{
``The score is the product of step-by-step probabilities, multiplied by a correctness gate that's 0 if any triple is fake---this makes the decoder faithful by construction.''
}

%% ------------------------------------------------------------
\subsection{5. TypeOracle Gates --- Semantic Pruning}

\whyt{
The graph constraint ensures the triple exists. The TypeOracle goes further: it checks whether the target entity's \emph{type} makes sense for the relation. This prunes bad candidates before the LLM even sees them, reducing the search space.
}

\stepboxbegin
\step{Setup: types in the KG.}
Every entity $e$ has a set of types $\text{rk}_e$ (e.g.\ Paris: \{City, Place\}). Every relation $r$ has a range type $\text{tr}_r$ (e.g.\ capital\_of: \{Country\}).

\step{The range gate.}
Given a relation $r$ and a candidate entity $e$, accept only if $e$'s type is compatible with $r$'s range:
\[
\text{range\_gate}(r, e) = \mathbb{1}[\text{tr}_r \subseteq \text{rk}_e]
\]
Example: capital\_of requires Country type. ``Paris'' has type City---gate rejects it. ``France'' has type Country---gate accepts it.

\step{The type gate.}
Given a head entity $s$ and relation $r$, compute what types the target could be:
\[
\text{type\_gate}(s, r) = \bigcup_{e \in \mathcal{N}(s,r)} \text{rk}_e
\]
This is the union of types of all actual neighbours of $s$ via $r$. It tells the LLM: ``the next entity must be one of these types.''

\step{How they work together.}
At each step:
\begin{enumerate}
  \item Compute $\text{type\_gate}(\text{head}, r)$ = what types are allowed.
  \item For each candidate entity $e$, check $\text{range\_gate}(r, e)$ = is $e$'s type in the allowed set?
  \item Only entities passing both gates are offered to the LLM.
\end{enumerate}

\step{Worked example.}
Head = ``Paris'', relation = ``capital\_of'':
\begin{itemize}
  \item type\_gate: neighbours of Paris via capital\_of = \{France\}. Types = \{Country\}.
  \item Candidate ``France'': range\_gate checks Country $\in$ \{Country\} $\to$ \textbf{accept}.
  \item Candidate ``Berlin'': range\_gate checks Country $\in$ \{City\} $\to$ \textbf{reject}.
\end{itemize}
\stepboxend

\oneliner{
``The TypeOracle is a semantic filter: it checks whether the candidate entity's type makes sense for the relation, pruning nonsensical paths before the LLM even considers them.''
}

%% ------------------------------------------------------------
\subsection{6. Complexity Analysis --- Why v2 Scales}

\whyt{
This is the core argument for why our approach works on real KGs. V1 hits a wall at chain length 3--4. V2 scales linearly.
}

\stepboxbegin
\step{V1: enumerate all paths.}
For each of $L$ steps, try all $E$ neighbours. Total paths: $E^L$. For each path, score it in $O(L)$ time. So:
\[
\text{V1} = O(E^L \cdot L)
\]
Example: $E = 100$, $L = 3$. That's $100^3 = 1{,}000{,}000$ paths. Feasible. But $L = 4$? $100{,}000{,}000$. Not feasible.

\step{V2: beam search.}
At each of $L$ steps, expand $B$ beams. Each beam expands to $d_{\text{avg}}$ candidates. Score each in $O(1)$ (we track the running product). So:
\[
\text{V2} = O(L \cdot B \cdot d_{\text{avg}})
\]
Example: $L = 4$, $B = 5$, $d_{\text{avg}} = 10$. That's $4 \times 5 \times 10 = 200$ operations. Always feasible.

\step{Where the savings come from.}
V1 is exponential because each step multiplies the number of paths. V2 is linear because we keep only the top-$B$ at each step---we never let the number of paths grow beyond $B$.

\step{What if $B$ is too small?}
We might miss the correct answer. This is the precision-recall tradeoff of beam search. Typical: $B = 3$--$10$.
\stepboxend

\oneliner{
``V1 tries all paths---exponential. V2 keeps only the best $B$ at each step---linear. That's why it scales to real knowledge graphs.''
}

%% ------------------------------------------------------------
\subsection{7. Beam Search --- How We Find the Best Path}

\whyt{
Beam search is how we go from ``enumerate everything'' (v1) to ``keep the best candidates'' (v2). It's the algorithm that makes v2 practical.
}

\stepboxbegin
\step{The idea.}
At each step:
\begin{enumerate}
  \item Take each beam (partial path).
  \item Expand it to all valid next entities (using the constrained LLM).
  \item Score each expansion.
  \item Keep only the top $B$ across all beams.
\end{enumerate}

\step{The dynamic trie (DoG contribution).}
Instead of pre-computing all valid paths, we build a trie on-the-fly as the LLM generates. Each beam has its own trie that grows as new triples are generated.

\step{Head pool (DoG contribution).}
Each beam tracks which entities have been mentioned so far ($\text{head}(i)$). This prevents cycles: you can't go back to an entity you already visited.

\step{Stopping condition.}
We stop when:
\begin{itemize}
  \item The LLM generates the $L$-th hop, or
  \item No valid expansions exist (dead end), or
  \item The LLM generates an end-of-sequence token.
\end{itemize}

\step{What we return.}
The beam with the highest score $S(\vec{c})$.
\stepboxend

\oneliner{
``Beam search keeps the best $B$ partial paths at each step, expanding only the most promising ones---it's like keeping a shortlist that gets refined at each hop.''
}

%% ------------------------------------------------------------
\subsection{8. Hits@k --- How We Measure Success}

\whyt{
We need a single number to say ``how good is our method?'' Hits@k tells us: out of all test questions, in what fraction was the correct answer in our top-$k$ predictions?
}

\stepboxbegin
\step{Definition.}
For a test question $q$, let $\text{rank}(q)$ be the position of the correct answer when we sort all candidate answers by score. Then:
\[
\text{Hits@}k = \frac{|\{q \in \mathcal{Q} : \text{rank}(q) \leq k\}|}{|\mathcal{Q}|}
\]
In words: ``what fraction of questions have the correct answer in the top $k$?''

\step{Hits@1.}
Is the \emph{very best} prediction correct? This is the strictest metric.

\step{Hits@3 / Hits@10.}
Is the correct answer \emph{somewhere} in the top 3 or top 10? More forgiving.

\step{What numbers are good?}
On WebQSP:
\begin{itemize}
  \item Baseline (no KG): $\sim$0.3
  \item GCR (ICML 2025): $\sim$0.45--0.55
  \item DCA-Trie v1: $\sim$0.50--0.60
  \item DCA-Trie v2: $\sim$0.55--0.65 (expected)
\end{itemize}

\step{Why Hits@1, not accuracy?}
Because we might generate multiple candidate answers. Hits@1 checks if the top one is right. Accuracy would require a single binary judgment per question.
\stepboxend

\oneliner{
``Hits@1 is the fraction of test questions where our top prediction is correct---it's the headline number for how well the method works.''
}

%% ------------------------------------------------------------
\subsection{9. Faithfulness Theorem --- Why the Output is Always Correct}

\whyt{
This is what makes DCA-Trie different from just asking an LLM a question. We can \emph{prove} the output is grounded in the KG.
}

\stepboxbegin
\step{What we claim.}
The constrained decoder only outputs chains $\vec{c}$ where every triple is in the KG:
\[
T(\vec{c}) = \prod_{i=1}^{L} \mathbb{1}[(\text{head}(i), r_i, c_{i+1}) \in \mathcal{G}] = 1
\]

\step{Proof by induction on $L$.}

\emph{Base case} ($L=1$): We only offer tokens from $\mathcal{N}(\text{head}(0))$. So $(c_1, r_1, c_2) \in \mathcal{G}$ by construction.

\emph{Inductive step}: Assume the decoder is faithful for chains of length $k$. For length $k+1$:
\begin{enumerate}
  \item By the inductive hypothesis, $\text{head}(k) = c_k$ and all previous triples are in $\mathcal{G}$.
  \item At step $k+1$, the decoder only offers tokens from $\mathcal{N}(\text{head}(k))$.
  \item So $(c_k, r_k, c_{k+1}) \in \mathcal{G}$.
  \item Therefore $T(\vec{c}) = 1$ for all steps.
\end{enumerate}

\step{What this means in practice.}
If we use the constrained decoder, we can never generate a hallucinated triple. Every fact in the output chain is verified against the KG.
\stepboxend

\oneliner{
``We prove by induction that the decoder can only generate chains where every triple is in the knowledge graph---hallucination is structurally impossible.''
}

\end{document}
